% ============================================================
% CHAPITRE 20 — Création de logos professionnels
% ============================================================

\chapter{Création de logos professionnels}

\noindent
\textit{Dans les chapitres précédents, nous avons appris à maîtriser Illustrator : dessin vectoriel, plume, manipulation avancée des formes, couleurs, typographie, lettrage. Nous savons désormais créer des formes précises, composer du texte, appliquer des couleurs. Toutes ces compétences convergent vers l'un des projets les plus emblématiques du design graphique : la création d'un logo.}

\vspace{1em}

\noindent
Le \terme{logo} est bien plus qu'une simple image : c'est la \emph{signature visuelle} d'une marque, d'une entreprise, d'une organisation, d'un produit. Il permet d'identifier, de reconnaître, de mémoriser. Il transmet des valeurs, une personnalité, une promesse. Un bon logo est \terme{simple}, \terme{mémorable}, \terme{intemporel}, \terme{versatile}, \terme{approprié}. Un mauvais logo est complexe, daté, illisible, incohérent.

\vspace{1em}

\noindent
La création d'un logo est un processus \terme{stratégique}, pas seulement esthétique. Elle repose sur une compréhension profonde du client, de son marché, de son public, de ses valeurs. Elle exige de la recherche, de l'analyse, de l'esquisse, de l'itération, de la précision. Un logo ne se crée pas en quelques minutes : il se construit, se travaille, se peaufine. C'est l'un des projets les plus exigeants du design graphique, mais aussi l'un des plus gratifiants.

\vspace{1em}

\noindent
Ce chapitre a pour objectif de te faire maîtriser la création de logos professionnels dans Illustrator. Nous commencerons par comprendre la fonction d'un logo. Nous étudierons ensuite les principales catégories de logos, l'analyse du brief, la recherche visuelle, l'esquisse de concepts, la construction d'un logo géométrique, la création d'un logo typographique, la conception d'un symbole vectoriel, l'évaluation de la lisibilité, le test des versions monochromes, la préparation des variantes, et l'organisation des fichiers sources et des exports. À la fin de ce chapitre, tu sauras créer des logos professionnels, de la conception à la livraison.

% ============================================================
\section{Comprendre la fonction d'un logo}
% ============================================================

\subsection{Qu'est-ce qu'un logo ?}

\begin{quote}
\textbf{Un logo est un signe graphique qui identifie et représente une marque, une entreprise, une organisation, un produit ou un service. Il combine généralement un symbole, une typographie, et une couleur, pour créer une identité visuelle reconnaissable et mémorable.}
\end{quote}

Le logo est la \terme{signature visuelle} d'une entité. Il apparaît sur tous les supports de communication : papeterie, site web, réseaux sociaux, packaging, publicité, signalétique. Il est le point de contact le plus fréquent entre une marque et son public.

\subsection{Les fonctions du logo}

\subsubsection{Identifier}

La première fonction du logo est d'\terme{identifier}. Il permet de reconnaître une marque parmi des milliers d'autres. Dans un marché saturé, un logo reconnaissable est un avantage concurrentiel majeur.

\subsubsection{Différencier}

Le logo permet de \terme{différencier} une marque de ses concurrents. Il affirme une personnalité, une identité, une position. Un logo distinctif se démarque, se retient, se préfère.

\subsubsection{Communiquer}

Le logo \terme{communique} des valeurs, une personnalité, une promesse. Il transmet des messages implicites : modernité, tradition, confiance, créativité, rigueur, chaleur. Un logo bien conçu dit quelque chose de la marque, sans avoir besoin de mots.

\subsubsection{Rassurer}

Le logo \terme{rassure}. Un logo professionnel inspire confiance, crédibilité, sérieux. À l'inverse, un logo approximatif suscite le doute, la méfiance.

\subsubsection{Fédérer}

Le logo \terme{fédère}. Il crée un sentiment d'appartenance chez les clients, les employés, les partenaires. Il devient un symbole identitaire, un point de ralliement.

\subsubsection{Marquer}

Le logo \terme{marque} les esprits. Un logo mémorable reste gravé dans la mémoire, associé à une expérience, à une émotion, à un souvenir.

\subsection{Les qualités d'un bon logo}

\subsubsection{La simplicité}

Un bon logo est \terme{simple} : peu de formes, peu de couleurs, peu de détails. La simplicité permet la reconnaissance immédiate, la mémorisation, la reproduction à toutes les tailles.

\subsubsection{La mémorisation}

Un bon logo est \terme{mémorable} : il marque, il reste, il se retient. Un logo complexe s'oublie ; un logo simple se retient.

\subsubsection{L'intemporalité}

Un bon logo est \terme{intemporel} : il ne suit pas les modes, il ne se démode pas. Il traverse les époques, les tendances, les technologies. Un logo daté doit être refait tous les dix ans ; un logo intemporel dure des décennies.

\subsubsection{La versatilité}

Un bon logo est \terme{versatile} : il fonctionne sur tous les supports, à toutes les tailles, en couleur et en noir et blanc, sur fond clair et sur fond sombre. Il s'adapte, il se décline, il reste reconnaissable.

\subsubsection{La pertinence}

Un bon logo est \terme{pertinent} : il correspond à la marque, à son marché, à son public. Il transmet les bonnes valeurs, la bonne personnalité, le bon message.

\subsubsection{L'appropriation}

Un bon logo est \terme{approprié} : il ne ressemble à aucun autre, il est unique, il est protégeable juridiquement. Un logo qui ressemble à un autre est un logo qui disparaît.

\subsection{Les erreurs à éviter}

\subsubsection{Trop complexe}

Un logo trop complexe est illisible à petite taille, difficile à mémoriser, impossible à reproduire sur certains supports.

\subsubsection{Trop tendance}

Un logo trop tendance se démode rapidement. Il faut privilégier l'intemporalité.

\subsubsection{Trop générique}

Un logo trop générique ne différencie pas la marque. Il faut chercher l'originalité, la distinction.

\subsubsection{Trop chargé}

Un logo trop chargé (trop de couleurs, trop de détails, trop d'effets) perd en lisibilité, en élégance, en impact.

\subsubsection{Copie}

Un logo copié est un logo illégal, qui nuit à la réputation de la marque et expose à des poursuites.

\subsection{Les bonnes pratiques}

\begin{itemize}
    \item \terme{Comprendre} la fonction du logo.
    \item \terme{Rechercher} la simplicité.
    \item \terme{Rechercher} la mémorisation.
    \item \terme{Rechercher} l'intemporalité.
    \item \terme{Rechercher} la versatilité.
    \item \terme{Rechercher} la pertinence.
    \item \terme{Éviter} la complexité, la tendance, la banalité.
    \item \terme{Vérifier} l'originalité.
\end{itemize}

\subsection*{Conclusion de la section}

Un \terme{logo} est un signe graphique qui identifie et représente une marque. Il remplit plusieurs fonctions : identifier, différencier, communiquer, rassurer, fédérer, marquer. Un bon logo est simple, mémorable, intemporel, versatile, pertinent, approprié. Il faut éviter les logos trop complexes, trop tendance, trop génériques, trop chargés, ou copiés.

Maintenant que nous comprenons la fonction d'un logo, nous allons, dans la section suivante, étudier les \terme{principales catégories de logos}.

% ============================================================
\section{Les principales catégories de logos}
% ============================================================

\subsection{Pourquoi connaître les catégories ?}

Connaître les catégories de logos permet de :

\begin{itemize}
    \item \terme{Choisir} le type de logo adapté à la marque.
    \item \terme{Comprendre} les avantages et inconvénients de chaque type.
    \item \terme{Communiquer} avec le client.
    \item \terme{S'inspirer} de références.
    \item \terme{Construire} un logo cohérent.
\end{itemize}

Les catégories ne sont pas rigides : de nombreux logos combinent plusieurs types.

\subsection{Les principales catégories}

\subsubsection{Le logo typographique (logotype)}

Le \terme{logotype} (ou \emph{wordmark}) est un logo composé uniquement de texte : le nom de la marque, écrit dans une typographie caractéristique.

\paragraph{Caractéristiques} :
\begin{itemize}
    \item \terme{Texte} comme élément principal.
    \item \terme{Typographie} distinctive.
    \item \terme{Pas de symbole}.
\end{itemize}

\paragraph{Avantages} :
\begin{itemize}
    \item \terme{Lisibilité} : le nom est immédiatement lisible.
    \item \terme{Mémorisation} : le nom est associé à une typographie.
    \item \terme{Simplicité} : facile à reproduire.
\end{itemize}

\paragraph{Inconvénients} :
\begin{itemize}
    \item \terme{Moins visuel} : pas de symbole mémorisable.
    \item \terme{Dépendant} du nom : un nom long est difficile.
\end{itemize}

\paragraph{Exemples} :
\begin{itemize}
    \item Coca-Cola.
    \item Google.
    \item FedEx.
    \item Sony.
    \item Vogue.
\end{itemize}

\subsubsection{Le logo symbolique (pictogramme)}

Le \terme{logo symbolique} est un logo composé uniquement d'un symbole ou d'un pictogramme, sans texte.

\paragraph{Caractéristiques} :
\begin{itemize}
    \item \terme{Symbole} comme élément principal.
    \item \terme{Pas de texte}.
    \item \terme{Forme} simplifiée, reconnaissable.
\end{itemize}

\paragraph{Avantages} :
\begin{itemize}
    \item \terme{Universalité} : le symbole est compris dans toutes les langues.
    \item \terme{Impact} : le symbole est visuel, mémorisable.
    \item \terme{Versatilité} : fonctionne à toutes les tailles.
\end{itemize}

\paragraph{Inconvénients} :
\begin{itemize}
    \item \terme{Ambiguïté} : le symbole peut être mal interprété.
    \item \terme{Nécessite} une reconnaissance préalable.
\end{itemize}

\paragraph{Exemples} :
\begin{itemize}
    \item Apple (pomme croquée).
    \item Nike (swoosh).
    \item Twitter (oiseau).
    \item Shell (coquillage).
    \item Mercedes (étoile à trois branches).
\end{itemize}

\subsubsection{Le logo mixte (combinaison)}

Le \terme{logo mixte} combine un symbole et un texte. C'est la catégorie la plus courante.

\paragraph{Caractéristiques} :
\begin{itemize}
    \item \terme{Symbole} et \terme{texte} associés.
    \item \terme{Plusieurs} configurations possibles (symbole au-dessus, à côté, intégré).
\end{itemize}

\paragraph{Avantages} :
\begin{itemize}
    \item \terme{Richesse} : combine les avantages du symbole et du texte.
    \item \terme{Flexibilité} : peut être décliné en symbole seul ou en texte seul.
    \item \terme{Reconnaissance} : le symbole et le texte se renforcent mutuellement.
\end{itemize}

\paragraph{Inconvénients} :
\begin{itemize}
    \item \terme{Complexité} : plus difficile à reproduire à petite taille.
    \item \terme{Nécessite} une bonne composition.
\end{itemize}

\paragraph{Exemples} :
\begin{itemize}
    \item Adidas.
    \item Lacoste.
    \item Puma.
    \item Starbucks.
    \item Airbnb.
\end{itemize}

\subsubsection{Le logo abstrait}

Le \terme{logo abstrait} est un logo composé d'une forme géométrique ou abstraite, sans lien direct avec le produit ou le service.

\paragraph{Caractéristiques} :
\begin{itemize}
    \item \terme{Forme} abstraite.
    \item \terme{Pas de symbolique} évidente.
    \item \terme{Interprétation} laissée au spectateur.
\end{itemize}

\paragraph{Avantages} :
\begin{itemize}
    \item \terme{Modernité} : les formes abstraites sont intemporelles.
    \item \terme{Distinction} : peu de logos abstraits existent.
    \item \terme{Flexibilité} : la forme peut évoquer plusieurs idées.
\end{itemize}

\paragraph{Inconvénients} :
\begin{itemize}
    \item \terme{Manque} de lien avec le produit.
    \item \terme{Nécessite} une reconnaissance préalable.
\end{itemize}

\paragraph{Exemples} :
\begin{itemize}
    \item Chanel (double C).
    \item BP (forme abstraite).
    \item Pepsi (cercle tricolore).
    \item Audi (quatre anneaux).
\end{itemize}

\subsubsection{Le logo emblème}

Le \terme{logo emblème} est un logo dans lequel le texte est intégré à une forme (blason, écusson, cercle).

\paragraph{Caractéristiques} :
\begin{itemize}
    \item \terme{Texte} intégré à une forme.
    \item \terme{Forme} fermée, délimitée.
    \item \terme{Style} souvent traditionnel, institutionnel.
\end{itemize}

\paragraph{Avantages} :
\begin{itemize}
    \item \terme{Impact} : forme fermée, mémorisable.
    \item \terme{Tradition} : évoque le sérieux, l'authenticité.
    \item \terme{Cohérence} : le texte et la forme ne font qu'un.
\end{itemize}

\paragraph{Inconvénients} :
\begin{itemize}
    \item \terme{Moins} versatile : difficile à décliner en symbole seul.
    \item \terme{Difficile} à lire à petite taille.
\end{itemize}

\paragraph{Exemples} :
\begin{itemize}
    \item Starbucks (sirène dans un cercle).
    \item Harley-Davidson (écusson).
    \item Harvard (blason).
    \item NFL (écusson).
\end{itemize}

\subsubsection{Le logo lettrage (lettering)}

Le \terme{logo lettrage} est un logo dessiné à la main, unique, avec des lettres personnalisées.

\paragraph{Caractéristiques} :
\begin{itemize}
    \item \terme{Lettres} dessinées à la main.
    \item \terme{Style} unique, personnalisé.
    \item \terme{Personnalité} forte.
\end{itemize}

\paragraph{Avantages} :
\begin{itemize}
    \item \terme{Originalité} : aucune autre marque n'a le même.
    \item \terme{Personnalité} : transmet une identité forte.
    \item \terme{Artisanat} : évoque le soin, la qualité.
\end{itemize}

\paragraph{Inconvénients} :
\begin{itemize}
    \item \terme{Coût} : plus long à créer.
    \item \terme{Moins} lisible à petite taille.
\end{itemize}

\paragraph{Exemples} :
\begin{itemize}
    \item Coca-Cola (lettrage original).
    \item Disney (lettrage signature).
    \item Instagram (lettrage moderne).
    \item Vans (lettrage).
\end{itemize}

\subsubsection{Le logo monogramme}

Le \terme{logo monogramme} est un logo composé des initiales de la marque, souvent entrelacées.

\paragraph{Caractéristiques} :
\begin{itemize}
    \item \terme{Initiales} comme élément principal.
    \item \terme{Forme} compacte, mémorisable.
    \item \terme{Style} souvent élégant, luxueux.
\end{itemize}

\paragraph{Avantages} :
\begin{itemize}
    \item \terme{Compact} : fonctionne à toutes les tailles.
    \item \terme{Élégant} : évoque le luxe, la tradition.
    \item \terme{Mémorisable} : les initiales sont simples.
\end{itemize}

\paragraph{Inconvénients} :
\begin{itemize}
    \item \terme{Nécessite} une reconnaissance préalable du nom.
    \item \terme{Limite} à deux ou trois initiales.
\end{itemize}

\paragraph{Exemples} :
\begin{itemize}
    \item Chanel (CC).
    \item Gucci (GG).
    \item Louis Vuitton (LV).
    \item Yves Saint Laurent (YSL).
\end{itemize}

\subsection{Tableau récapitulatif des catégories}

\begin{table}[h!]
\centering
\begin{tabular}{@{}p{3cm}p{4cm}p{4cm}p{3cm}@{}}
\toprule
\textbf{Catégorie} & \textbf{Composition} & \textbf{Avantages} & \textbf{Exemples} \\
\midrule
Typographique & Texte seul & Lisibilité, mémorisation & Coca-Cola, Google \\
Symbolique & Symbole seul & Universalité, impact & Apple, Nike \\
Mixte & Symbole + texte & Richesse, flexibilité & Adidas, Starbucks \\
Abstrait & Forme abstraite & Modernité, distinction & Chanel, Audi \\
Emblème & Texte dans une forme & Impact, tradition & Starbucks, NFL \\
Lettrage & Lettres dessinées & Originalité, personnalité & Disney, Vans \\
Monogramme & Initiales & Compacité, élégance & Gucci, LV \\
\bottomrule
\end{tabular}
\caption{Les principales catégories de logos.}
\end{table}

\subsection{Choisir la bonne catégorie}

Le choix de la catégorie dépend de :

\begin{itemize}
    \item Du \terme{nom} de la marque (long ou court).
    \item Du \terme{secteur} (luxe, tech, alimentaire, etc.).
    \item Du \terme{public} (âge, milieu, culture).
    \item De la \terme{personnalité} (moderne, traditionnelle, créative).
    \item De la \terme{stratégie} (notoriété, distinction, confiance).
\end{itemize}

\subsection{Les bonnes pratiques}

\begin{itemize}
    \item \terme{Choisir} la catégorie adaptée à la marque.
    \item \terme{Combiner} plusieurs catégories si nécessaire.
    \item \terme{Vérifier} la lisibilité à toutes les tailles.
    \item \terme{Tester} en couleur et en noir et blanc.
    \item \terme{S'inspirer} de références.
    \item \terme{Éviter} de copier.
\end{itemize}

\subsection*{Conclusion de la section}

Il existe plusieurs catégories de logos : typographique (logotype), symbolique (pictogramme), mixte (combinaison), abstrait, emblème, lettrage, monogramme. Chaque catégorie a ses avantages et inconvénients. Le choix dépend du nom, du secteur, du public, de la personnalité et de la stratégie de la marque. Il est essentiel de choisir la bonne catégorie, de combiner si nécessaire, de vérifier la lisibilité.

Maintenant que nous connaissons les catégories, nous allons, dans la section suivante, apprendre à \terme{analyser un brief de création}.

% ============================================================
\section{Analyser un brief de création}
% ============================================================

\subsection{Qu'est-ce qu'un brief ?}

\begin{quote}
\textbf{Un brief (ou \emph{cahier des charges}) est un document qui définit le projet, les objectifs, le public, les contraintes, les références. Il sert de base à la création du logo.}
\end{quote}

Le brief est le point de départ de tout projet de logo. Il permet de comprendre le besoin, de cadrer le travail, d'éviter les malentendus. Un brief mal compris entraîne un travail inutile.

\subsection{Le contenu du brief}

\subsubsection{Le contexte}

\begin{itemize}
    \item \terme{Qui} est le client ?
    \item \terme{Que} fait-il ?
    \item \terme{Quel} est son marché ?
    \item \terme{Quelle} est son histoire ?
    \item \terme{Quels} sont ses concurrents ?
\end{itemize}

\subsubsection{Le projet}

\begin{itemize}
    \item \terme{Quel} est le projet ? (création, refonte)
    \item \terme{Pourquoi} ce projet ? (nouveau produit, changement de positionnement)
    \item \terme{Quand} doit-il être livré ?
    \item \terme{Quel} est le budget ?
\end{itemize}

\subsubsection{Les objectifs}

\begin{itemize}
    \item \terme{Quel} est l'objectif du logo ? (notoriété, distinction, confiance)
    \item \terme{Quel} message doit-il transmettre ?
    \item \terme{Quel} effet doit-il produire ?
\end{itemize}

\subsubsection{Le public}

\begin{itemize}
    \item \terme{Qui} est le public cible ? (âge, milieu, culture)
    \item \terme{Que} veut-il ?
    \item \terme{Quels} sont ses codes ?
    \item \terme{Comment} perçoit-il la marque ?
\end{itemize}

\subsubsection{Les contraintes}

\begin{itemize}
    \item \terme{Format} : quels supports ?
    \item \terme{Technique} : quelles contraintes (impression, web, textile) ?
    \item \terme{Juridique} : quelles contraintes (marques existantes, droits) ?
    \item \terme{Graphique} : y a-t-il une charte existante ?
\end{itemize}

\subsubsection{Les références}

\begin{itemize}
    \item \terme{Ce} que le client aime.
    \item \terme{Ce} qu'il n'aime pas.
    \item \terme{Les} logos qu'il admire.
    \item \terme{Les} logos qu'il rejette.
\end{itemize}

\subsection{Analyser le brief}

\subsubsection{Étape 1 : Lire attentivement}

Lire le brief plusieurs fois, noter les points importants, les questions à poser.

\subsubsection{Étape 2 : Poser des questions}

Le brief n'est pas toujours complet. Poser des questions pour clarifier :

\begin{itemize}
    \item \terme{Le} positionnement.
    \item \terme{Les} valeurs.
    \item \terme{Les} contraintes.
    \item \terme{Les} préférences.
\end{itemize}

\subsubsection{Étape 3 : Reformuler}

Reformuler le brief avec ses propres mots, pour vérifier la compréhension.

\subsubsection{Étape 4 : Valider}

Faire valider la reformulation par le client.

\subsection{Les questions à poser}

\subsubsection{Sur la marque}

\begin{itemize}
    \item Quelles sont les valeurs de la marque ?
    \item Quelle est sa personnalité ?
    \item Quelle est son histoire ?
    \item Quels sont ses concurrents ?
\end{itemize}

\subsubsection{Sur le public}

\begin{itemize}
    \item Qui est le public cible ?
    \item Quel est son âge, son milieu, sa culture ?
    \item Quels sont ses codes ?
    \item Comment perçoit-il la marque ?
\end{itemize}

\subsubsection{Sur le projet}

\begin{itemize}
    \item Pourquoi ce projet ?
    \item Quand doit-il être livré ?
    \item Quel est le budget ?
    \item Y a-t-il des contraintes ?
\end{itemize}

\subsubsection{Sur les références}

\begin{itemize}
    \item Quels logos aimez-vous ?
    \item Quels logos n'aimez-vous pas ?
    \item Pourquoi ?
\end{itemize}

\subsection{Les bonnes pratiques}

\begin{itemize}
    \item \terme{Lire} attentivement le brief.
    \item \terme{Poser} des questions.
    \item \terme{Reformuler} pour vérifier.
    \item \terme{Valider} avec le client.
    \item \terme{Noter} les points clés.
    \item \terme{Documenter} les échanges.
\end{itemize}

\subsection*{Conclusion de la section}

Un \terme{brief} est un document qui définit le projet, les objectifs, le public, les contraintes, les références. Il contient le contexte, le projet, les objectifs, le public, les contraintes, les références. L'analyse suit plusieurs étapes : lire, poser des questions, reformuler, valider. Il est essentiel de lire attentivement, de poser des questions, de reformuler, de valider avec le client.

Maintenant que nous savons analyser un brief, nous allons, dans la section suivante, apprendre à \terme{effectuer des recherches visuelles}.

% ============================================================
\section{Effectuer des recherches visuelles}
% ============================================================

\subsection{Pourquoi rechercher des références ?}

La \terme{recherche visuelle} est essentielle pour :

\begin{itemize}
    \item \terme{Comprendre} le marché et la concurrence.
    \item \terme{S'inspirer} de références.
    \item \terme{Éviter} de réinventer la roue.
    \item \terme{Développer} son vocabulaire visuel.
    \item \terme{Communiquer} avec le client.
    \item \terme{Se positionner} par rapport à l'existant.
\end{itemize}

\subsection{Les sources de références}

\subsubsection{Les plateformes en ligne}

\begin{itemize}
    \item \terme{Behance} : portfolios de designers.
    \item \terme{Dribbble} : projets de designers.
    \item \terme{Pinterest} : épingles visuelles.
    \item \terme{Logo Lounge} : bibliothèque de logos.
    \item \terme{Brand New} : critiques de logos.
    \item \terme{Logopond} : inspiration de logos.
\end{itemize}

\subsubsection{Les livres}

\begin{itemize}
    \item \terme{Livres de logos} : recueils de logos.
    \item \terme{Livres de branding} : ouvrages sur les marques.
    \item \terme{Livres d'histoire du design} : références historiques.
\end{itemize}

\subsubsection{Le quotidien}

\begin{itemize}
    \item \terme{Rue} : logos, enseignes, panneaux.
    \item \terme{Magasins} : produits, packagings.
    \item \terme{Publicités} : campagnes, affiches.
    \item \terme{Applications} : interfaces, icônes.
\end{itemize}

\subsubsection{Les autres arts}

\begin{itemize}
    \item \terme{Peinture} : compositions, couleurs.
    \item \terme{Architecture} : formes, structures.
    \item \terme{Typographie} : lettres, styles.
    \item \terme{Illustration} : styles, techniques.
\end{itemize}

\subsection{Méthode de recherche}

\subsubsection{Étape 1 : Définir les mots-clés}

\begin{itemize}
    \item Mots-clés en lien avec le projet (par exemple, « logo tech », « logo luxe »).
    \item Mots-clés en lien avec le style (par exemple, « minimaliste », « géométrique »).
    \item Mots-clés en lien avec le secteur (par exemple, « café », « banque »).
\end{itemize}

\subsubsection{Étape 2 : Explorer les sources}

Explorer les différentes sources : plateformes, livres, quotidien, autres arts.

\subsubsection{Étape 3 : Sélectionner les références}

Sélectionner les références les plus pertinentes, en fonction de :

\begin{itemize}
    \item De leur \terme{qualité} graphique.
    \item De leur \terme{pertinence} par rapport au projet.
    \item De leur \terme{diversité} (ne pas se limiter à un seul style).
    \item De leur \terme{originalité}.
\end{itemize}

\subsubsection{Étape 4 : Analyser les références}

Analyser les références :

\begin{itemize}
    \item \terme{Quel} type de logo ?
    \item \terme{Quelle} typographie ?
    \item \terme{Quelle} couleur ?
    \item \terme{Quelle} composition ?
    \item \terme{Quelle} personnalité ?
    \item \terme{Quel} public ?
\end{itemize}

\subsubsection{Étape 5 : Constituer un moodboard}

Créer un moodboard (planche de références) pour organiser les références et communiquer avec le client.

\subsection{Les bonnes pratiques}

\begin{itemize}
    \item \terme{Définir} les mots-clés.
    \item \terme{Explorer} plusieurs sources.
    \item \terme{Sélectionner} les références pertinentes.
    \item \terme{Analyser} les références.
    \item \terme{Constituer} un moodboard.
    \item \terme{Éviter} de copier.
    \item \terme{Documenter} les sources.
\end{itemize}

\subsection*{Conclusion de la section}

La \terme{recherche visuelle} est essentielle pour comprendre, s'inspirer, se positionner. Les sources sont variées : plateformes en ligne, livres, quotidien, autres arts. La méthode suit plusieurs étapes : définir les mots-clés, explorer, sélectionner, analyser, constituer un moodboard. Il est essentiel de définir les mots-clés, d'explorer plusieurs sources, d'analyser, d'éviter de copier.

Maintenant que nous savons effectuer des recherches, nous allons, dans la section suivante, apprendre à \terme{esquisser plusieurs concepts}.

% ============================================================
\section{Esquisser plusieurs concepts}
% ============================================================

\subsection{Pourquoi esquisser ?}

L'\terme{esquisse} est une étape essentielle de la création d'un logo. Elle permet de :

\begin{itemize}
    \item \terme{Explorer} plusieurs pistes.
    \item \terme{Tester} des idées rapidement.
    \item \terme{Éviter} de se fixer sur une seule idée.
    \item \terme{Comparer} les concepts.
    \item \terme{Sélectionner} la meilleure piste.
    \item \terme{Présenter} plusieurs options au client.
\end{itemize}

L'esquisse se fait généralement sur papier, avant de passer à l'ordinateur. Le papier permet d'explorer rapidement, sans contrainte technique.

\subsection{Les principes de l'esquisse}

\subsubsection{La quantité}

Il faut \terme{beaucoup} esquisser. Plus on esquisse, plus on a de chances de trouver une bonne idée. Il ne faut pas se contenter de deux ou trois croquis : il en faut dix, vingt, trente.

\subsubsection{La rapidité}

Il faut esquisser \terme{rapidement}. L'esquisse n'a pas besoin d'être parfaite : elle doit juste capturer l'idée. On peut utiliser un crayon, un feutre, une tablette graphique.

\subsubsection{La diversité}

Il faut esquisser \terme{diverses} pistes : différents types de logos, différents styles, différentes idées. Il ne faut pas se limiter à une seule direction.

\subsubsection{La simplicité}

Il faut esquisser \terme{simplement}. Un bon logo est simple : l'esquisse doit l'être aussi. On évite les détails, les effets, les fioritures.

\subsection{Les étapes de l'esquisse}

\subsubsection{Étape 1 : Se documenter}

Relire le brief, regarder les références, s'imprégner du projet.

\subsubsection{Étape 2 : Brainstormer}

Générer des idées : mots-clés, concepts, symboles, formes, typographies. Ne pas se censurer.

\subsubsection{Étape 3 : Esquisser rapidement}

Dessiner rapidement plusieurs pistes : 10, 20, 30 croquis. Ne pas s'attarder sur chaque croquis.

\subsubsection{Étape 4 : Sélectionner}

Sélectionner les 3 à 5 pistes les plus prometteuses. Les développer un peu plus.

\subsubsection{Étape 5 : Raffiner}

Raffiner les pistes sélectionnées : préciser les formes, les proportions, les typographies.

\subsubsection{Étape 6 : Présenter}

Présenter les pistes au client, avec explication des choix.

\subsection{Les techniques d'esquisse}

\subsubsection{Le croquis papier}

Le \terme{croquis papier} est la technique la plus classique. Il permet d'explorer rapidement, sans contrainte technique.

\subsubsection{Le croquis numérique}

Le \terme{croquis numérique} se fait sur tablette graphique. Il permet de profiter des outils numériques (annulation, calques, couleurs).

\subsubsection{Le mind mapping}

Le \terme{mind mapping} consiste à créer une carte mentale autour du projet : mots-clés, concepts, associations. Il aide à générer des idées.

\subsubsection{Le brainstorming visuel}

Le \terme{brainstorming visuel} consiste à associer des images, des formes, des mots pour générer des idées.

\subsection{Les bonnes pratiques}

\begin{itemize}
    \item \terme{Esquisser} beaucoup.
    \item \terme{Esquisser} rapidement.
    \item \terme{Esquisser} diverses pistes.
    \item \terme{Esquisser} simplement.
    \item \terme{Ne pas se censurer}.
    \item \terme{Sélectionner} les meilleures pistes.
    \item \terme{Raffiner} avant de présenter.
    \item \terme{Documenter} les esquisses.
\end{itemize}

\subsection{Les erreurs courantes}

\subsubsection{Se fixer trop vite}

Se fixer sur une idée dès le début, sans explorer.

\subsubsection{Esquisser trop peu}

Ne faire que deux ou trois croquis, sans explorer.

\subsubsection{Esquisser trop complexe}

Vouloir tout de suite faire un logo parfait, avec des détails.

\subsubsection{Ne pas présenter assez d'options}

Présenter une seule piste au client, sans alternative.

\subsection*{Conclusion de la section}

L'\terme{esquisse} est essentielle pour explorer plusieurs pistes, tester des idées, comparer les concepts. Elle repose sur la quantité, la rapidité, la diversité, la simplicité. La méthode suit plusieurs étapes : se documenter, brainstormer, esquisser, sélectionner, raffiner, présenter. Il est essentiel d'esquisser beaucoup, rapidement, diverses pistes, sans se censurer.

Maintenant que nous savons esquisser, nous allons, dans la section suivante, apprendre à \terme{construire un logo géométrique}.

% ============================================================
\section{Construire un logo géométrique}
% ============================================================

\subsection{Qu'est-ce qu'un logo géométrique ?}

\begin{quote}
\textbf{Un logo géométrique est un logo composé de formes géométriques simples (cercles, carrés, triangles, lignes), assemblées pour créer un symbole.}
\end{quote}

Les logos géométriques sont utilisés pour :

\begin{itemize}
    \item Les \terme{marques modernes} : tech, finance, conseil.
    \item Les \terme{marques abstraites} : pas de symbole figuratif.
    \item Les \terme{identités minimalistes} : épurées, élégantes.
    \item Les \terme{logos intemporels} : formes simples, durables.
\end{itemize}

\subsection{Les principes du logo géométrique}

\subsubsection{La simplicité}

Un logo géométrique est \terme{simple} : peu de formes, peu de couleurs.

\subsubsection{La précision}

Un logo géométrique est \terme{précis} : les formes sont alignées, les proportions sont justes, les angles sont exacts.

\subsubsection{L'équilibre}

Un logo géométrique est \terme{équilibré} : les masses sont réparties, les formes sont harmonieuses.

\subsubsection{La cohérence}

Un logo géométrique est \terme{cohérent} : les formes partagent un style, une épaisseur, une logique.

\subsection{Construire un logo géométrique}

\subsubsection{Étape 1 : Choisir les formes de base}

\begin{itemize}
    \item \terme{Cercle} : harmonie, unité, mouvement.
    \item \terme{Carré} : stabilité, rigueur, équilibre.
    \item \terme{Triangle} : dynamisme, ascension, énergie.
    \item \terme{Ligne} : direction, connexion, mouvement.
\end{itemize}

\subsubsection{Étape 2 : Composer les formes}

\begin{itemize}
    \item \terme{Superposer} les formes.
    \item \terme{Combiner} les formes (Pathfinder).
    \item \terme{Aligner} les formes.
    \item \terme{Équilibrer} les formes.
\end{itemize}

\subsubsection{Étape 3 : Ajuster les proportions}

\begin{itemize}
    \item \terme{Vérifier} les proportions (nombre d'or, grille).
    \item \terme{Ajuster} les tailles.
    \item \terme{Ajuster} les positions.
\end{itemize}

\subsubsection{Étape 4 : Appliquer les couleurs}

\begin{itemize}
    \item \terme{Choisir} une palette cohérente.
    \item \terme{Appliquer} les couleurs.
    \item \terme{Vérifier} le contraste.
\end{itemize}

\subsubsection{Étape 5 : Vérifier}

\begin{itemize}
    \item \terme{Vérifier} la lisibilité à toutes les tailles.
    \item \terme{Vérifier} en noir et blanc.
    \item \terme{Vérifier} l'équilibre.
\end{itemize}

\subsection{Les techniques de construction}

\subsubsection{La grille}

La \terme{grille} est un outil essentiel pour construire un logo géométrique. Elle permet d'aligner les formes, de vérifier les proportions, de créer un logo cohérent.

\subsubsection{Les guides}

Les \terme{guides} (repères) permettent d'aligner les formes avec précision.

\subsubsection{Les formes de base}

Les \terme{formes de base} (cercle, carré, triangle) permettent de construire des logos géométriques simples et efficaces.

\subsubsection{Le Pathfinder}

Le \terme{Pathfinder} permet de combiner les formes (ajouter, soustraire, intersecter).

\subsection{Les exemples}

\subsubsection{Logo avec un cercle}

\begin{enumerate}
    \item Créer un cercle.
    \item Créer un carré.
    \item Soustraire le carré du cercle.
    \item Appliquer une couleur.
    \item Vérifier le résultat.
\end{enumerate}

\subsubsection{Logo avec un triangle}

\begin{enumerate}
    \item Créer un triangle.
    \item Créer un cercle.
    \item Combiner les formes.
    \item Appliquer une couleur.
    \item Vérifier le résultat.
\end{enumerate}

\subsubsection{Logo avec des lignes}

\begin{enumerate}
    \item Créer des lignes.
    \item Aligner les lignes.
    \item Combiner les lignes.
    \item Appliquer une couleur.
    \item Vérifier le résultat.
\end{enumerate}

\subsection{Les bonnes pratiques}

\begin{itemize}
    \item \terme{Utiliser} des formes simples.
    \item \terme{Aligner} les formes.
    \item \terme{Vérifier} les proportions.
    \item \terme{Utiliser} une grille.
    \item \terme{Appliquer} une palette cohérente.
    \item \terme{Vérifier} la lisibilité.
    \item \terme{Vérifier} en noir et blanc.
\end{itemize}

\subsection*{Conclusion de la section}

Un \terme{logo géométrique} est composé de formes simples : cercles, carrés, triangles, lignes. Il repose sur la simplicité, la précision, l'équilibre, la cohérence. La construction suit plusieurs étapes : choisir les formes, composer, ajuster, appliquer les couleurs, vérifier. Il est essentiel d'utiliser des formes simples, d'aligner, de vérifier les proportions, de tester la lisibilité.

Maintenant que nous savons construire un logo géométrique, nous allons, dans la section suivante, apprendre à \terme{créer un logo typographique}.

% ============================================================
\section{Créer un logo typographique}
% ============================================================

\subsection{Qu'est-ce qu'un logo typographique ?}

\begin{quote}
\textbf{Un logo typographique (ou logotype) est un logo composé uniquement de texte : le nom de la marque, écrit dans une typographie caractéristique.}
\end{quote}

Les logos typographiques sont utilisés pour :

\begin{itemize}
    \item Les \terme{marques au nom fort} : le nom suffit à identifier.
    \item Les \terme{marques modernes} : simplicité, élégance.
    \item Les \terme{marques de luxe} : typographie raffinée.
    \item Les \terme{marques intemporelles} : le nom reste, la typographie aussi.
\end{itemize}

\subsection{Les principes du logo typographique}

\subsubsection{La typographie}

La \terme{typographie} est l'élément principal. Elle doit être :

\begin{itemize}
    \item \terme{Lisible} : à toutes les tailles.
    \item \terme{Distinctive} : reconnaissable entre toutes.
    \item \terme{Adaptée} : au message, au public, au secteur.
    \item \terme{Cohérente} : avec l'identité de la marque.
\end{itemize}

\subsubsection{La composition}

La \terme{composition} du texte doit être :

\begin{itemize}
    \item \terme{Équilibrée} : les lettres sont bien réparties.
    \item \terme{Alignée} : les lettres sont alignées.
    \item \terme{Espacée} : les espacements sont ajustés.
    \item \terme{Harmonieuse} : l'ensemble forme un tout.
\end{itemize}

\subsubsection{La personnalisation}

La \terme{personnalisation} permet de rendre le logo unique :

\begin{itemize}
    \item \terme{Modifier} les lettres (allonger, élargir, courber).
    \item \terme{Ajouter} des fioritures.
    \item \terme{Combiner} les lettres.
    \item \terme{Créer} des lettres uniques.
\end{itemize}

\subsection{Créer un logo typographique}

\subsubsection{Étape 1 : Choisir la police}

\begin{itemize}
    \item \terme{Analyser} le brief, le public, le secteur.
    \item \terme{Choisir} une police adaptée (serif, sans-serif, script).
    \item \terme{Tester} plusieurs polices.
    \item \terme{Sélectionner} la plus pertinente.
\end{itemize}

\subsubsection{Étape 2 : Composer le texte}

\begin{itemize}
    \item \terme{Saisir} le texte.
    \item \terme{Ajuster} la taille.
    \item \terme{Ajuster} l'interlignage.
    \item \terme{Ajuster} le crénage.
    \item \terme{Ajuster} l'approche.
\end{itemize}

\subsubsection{Étape 3 : Personnaliser}

\begin{itemize}
    \item \terme{Vectoriser} le texte.
    \item \terme{Modifier} les lettres.
    \item \terme{Ajouter} des fioritures.
    \item \terme{Combiner} les lettres.
\end{itemize}

\subsubsection{Étape 4 : Appliquer les couleurs}

\begin{itemize}
    \item \terme{Choisir} une palette cohérente.
    \item \terme{Appliquer} les couleurs.
    \item \terme{Vérifier} le contraste.
\end{itemize}

\subsubsection{Étape 5 : Vérifier}

\begin{itemize}
    \item \terme{Vérifier} la lisibilité à toutes les tailles.
    \item \terme{Vérifier} en noir et blanc.
    \item \terme{Vérifier} l'équilibre.
    \item \terme{Vérifier} la cohérence.
\end{itemize}

\subsection{Les techniques de personnalisation}

\subsubsection{Modifier les lettres}

\begin{itemize}
    \item \terme{Allonger} une lettre.
    \item \terme{Élargir} une lettre.
    \item \terme{Incliner} une lettre.
    \item \terme{Courber} une lettre.
\end{itemize}

\subsubsection{Ajouter des fioritures}

\begin{itemize}
    \item \terme{Ajouter} des extensions.
    \item \terme{Ajouter} des ornements.
    \item \terme{Ajouter} des ligatures.
\end{itemize}

\subsubsection{Combiner les lettres}

\begin{itemize}
    \item \terme{Lier} deux lettres (ligature).
    \item \terme{Superposer} deux lettres.
    \item \terme{Entrelacer} deux lettres.
\end{itemize}

\subsection{Les exemples}

\subsubsection{Logo typographique simple}

\begin{enumerate}
    \item Choisir une police.
    \item Saisir le texte.
    \item Ajuster la taille, l'interlignage, le crénage.
    \item Appliquer une couleur.
    \item Vérifier le résultat.
\end{enumerate}

\subsubsection{Logo typographique personnalisé}

\begin{enumerate}
    \item Vectoriser le texte.
    \item Modifier les lettres.
    \item Ajouter des fioritures.
    \item Appliquer une couleur.
    \item Vérifier le résultat.
\end{enumerate}

\subsection{Les bonnes pratiques}

\begin{itemize}
    \item \terme{Choisir} une police adaptée.
    \item \terme{Ajuster} les espacements.
    \item \terme{Personnaliser} les lettres.
    \item \terme{Vérifier} la lisibilité.
    \item \terme{Vérifier} en noir et blanc.
    \item \terme{Conserver} une version modifiable.
    \item \terme{Éviter} les polices trop communes.
\end{itemize}

\subsection*{Conclusion de la section}

Un \terme{logo typographique} est composé uniquement de texte. Il repose sur la typographie, la composition, la personnalisation. La construction suit plusieurs étapes : choisir la police, composer, personnaliser, appliquer les couleurs, vérifier. Il est essentiel de choisir une police adaptée, d'ajuster les espacements, de personnaliser, de vérifier la lisibilité.

Maintenant que nous savons créer un logo typographique, nous allons, dans la section suivante, apprendre à \terme{créer un symbole vectoriel}.

% ============================================================
\section{Créer un symbole vectoriel}
% ============================================================

\subsection{Qu'est-ce qu'un symbole vectoriel ?}

\begin{quote}
\textbf{Un symbole vectoriel est une forme graphique simple, reconnaissable, qui représente la marque. Il peut être figuratif (objet, animal, personne) ou abstrait (forme géométrique).}
\end{quote}

Les symboles vectoriels sont utilisés pour :

\begin{itemize}
    \item Les \terme{logos symboliques} : symbole seul.
    \item Les \terme{logos mixtes} : symbole + texte.
    \item Les \terme{icônes} : versions simplifiées du logo.
    \item Les \terme{applications} : favicon, avatar.
\end{itemize}

\subsection{Les principes du symbole vectoriel}

\subsubsection{La simplicité}

Un symbole vectoriel est \terme{simple} : peu de formes, peu de détails, peu de couleurs.

\subsubsection{La reconnaissance}

Un symbole vectoriel est \terme{reconnaissable} : il évoque immédiatement la marque, le produit, le service.

\subsubsection{La mémorisation}

Un symbole vectoriel est \terme{mémorable} : il marque, il reste, il se retient.

\subsubsection{La versatilité}

Un symbole vectoriel est \terme{versatile} : il fonctionne à toutes les tailles, en couleur et en noir et blanc.

\subsection{Créer un symbole vectoriel}

\subsubsection{Étape 1 : Concevoir}

\begin{itemize}
    \item \terme{Analyser} le brief.
    \item \terme{Chercher} des idées (brainstorming).
    \item \terme{Esquisser} plusieurs pistes.
    \item \terme{Sélectionner} la meilleure piste.
\end{itemize}

\subsubsection{Étape 2 : Dessiner les formes de base}

\begin{itemize}
    \item \terme{Dessiner} avec la Plume.
    \item \terme{Utiliser} des formes géométriques.
    \item \terme{Combiner} les formes.
\end{itemize}

\subsubsection{Étape 3 : Simplifier}

\begin{itemize}
    \item \terme{Supprimer} les détails inutiles.
    \item \terme{Simplifier} les tracés.
    \item \terme{Réduire} le nombre de points.
\end{itemize}

\subsubsection{Étape 4 : Appliquer les couleurs}

\begin{itemize}
    \item \terme{Choisir} une palette cohérente.
    \item \terme{Appliquer} les couleurs.
    \item \terme{Vérifier} le contraste.
\end{itemize}

\subsubsection{Étape 5 : Vérifier}

\begin{itemize}
    \item \terme{Vérifier} la lisibilité à toutes les tailles.
    \item \terme{Vérifier} en noir et blanc.
    \item \terme{Vérifier} la reconnaissance.
    \item \terme{Vérifier} la cohérence.
\end{itemize}

\subsection{Les techniques}

\subsubsection{Le dessin avec la Plume}

\begin{itemize}
    \item \terme{Tracer} les contours.
    \item \terme{Ajuster} les courbes.
    \item \terme{Simplifier} les tracés.
\end{itemize}

\subsubsection{La construction géométrique}

\begin{itemize}
    \item \terme{Utiliser} des formes géométriques.
    \item \terme{Combiner} les formes.
    \item \terme{Aligner} les formes.
\end{itemize}

\subsubsection{La simplification}

\begin{itemize}
    \item \terme{Supprimer} les points inutiles.
    \item \terme{Réduire} le nombre de formes.
    \item \terme{Réduire} le nombre de couleurs.
\end{itemize}

\subsection{Les exemples}

\subsubsection{Symbole figuratif}

\begin{enumerate}
    \item Choisir un objet (animal, plante, objet).
    \item Esquisser plusieurs pistes.
    \item Dessiner avec la Plume.
    \item Simplifier.
    \item Appliquer les couleurs.
    \item Vérifier.
\end{enumerate}

\subsubsection{Symbole abstrait}

\begin{enumerate}
    \item Choisir des formes géométriques.
    \item Combiner les formes.
    \item Ajuster les proportions.
    \item Appliquer les couleurs.
    \item Vérifier.
\end{enumerate}

\subsection{Les bonnes pratiques}

\begin{itemize}
    \item \terme{Concevoir} avant de dessiner.
    \item \terme{Dessiner} avec la Plume.
    \item \terme{Simplifier} au maximum.
    \item \terme{Vérifier} la reconnaissance.
    \item \terme{Vérifier} la lisibilité.
    \item \terme{Vérifier} en noir et blanc.
    \item \terme{Tester} à toutes les tailles.
\end{itemize}

\subsection*{Conclusion de la section}

Un \terme{symbole vectoriel} est une forme graphique simple qui représente la marque. Il repose sur la simplicité, la reconnaissance, la mémorisation, la versatilité. La construction suit plusieurs étapes : concevoir, dessiner, simplifier, appliquer les couleurs, vérifier. Il est essentiel de concevoir avant de dessiner, de simplifier, de vérifier la reconnaissance et la lisibilité.

Maintenant que nous savons créer un symbole vectoriel, nous allons, dans la section suivante, apprendre à \terme{évaluer la lisibilité à différentes tailles}.

% ============================================================
\section{Évaluer la lisibilité à différentes tailles}
% ============================================================

\subsection{Pourquoi évaluer la lisibilité ?}

La \terme{lisibilité} est essentielle pour un logo. Un logo doit être lisible :

\begin{itemize}
    \item \terme{À} toutes les tailles.
    \item \terme{Sur} tous les supports.
    \item \terme{Dans} tous les contextes.
\end{itemize}

Un logo illisible à petite taille (favicon, avatar, carte de visite) est un logo raté. Un logo illisible à grande taille (affiche, panneau) est un logo raté. Il faut donc tester la lisibilité à différentes tailles.

\subsection{Les tailles à tester}

\subsubsection{Très petite taille}

\begin{itemize}
    \item \terme{Favicon} : 16 × 16 px.
    \item \terme{Avatar} : 32 × 32 px.
    \item \terme{Carte de visite} : 85 × 55 mm.
\end{itemize}

\subsubsection{Petite taille}

\begin{itemize}
    \item \terme{En-tête de site} : 100 × 100 px.
    \item \terme{Réseaux sociaux} : 400 × 400 px.
    \item \terme{Signature email} : 150 × 50 px.
\end{itemize}

\subsubsection{Taille moyenne}

\begin{itemize}
    \item \terme{Affiche A3} : 297 × 420 mm.
    \item \terme{Papeterie} : A4.
    \item \terme{Bannière web} : 728 × 90 px.
\end{itemize}

\subsubsection{Grande taille}

\begin{itemize}
    \item \terme{Affiche A0} : 841 × 1189 mm.
    \item \terme{Panneau} : 4 × 3 m.
    \item \terme{Enseigne} : variable.
\end{itemize}

\subsection{Comment évaluer la lisibilité ?}

\subsubsection{Test à l'écran}

\begin{enumerate}
    \item Afficher le logo à différentes tailles à l'écran.
    \item Zoomer et dézoomer.
    \item Vérifier la lisibilité.
\end{enumerate}

\subsubsection{Test sur papier}

\begin{enumerate}
    \item Imprimer le logo à différentes tailles.
    \item Vérifier la lisibilité.
    \item Vérifier la qualité.
\end{enumerate}

\subsubsection{Test en contexte}

\begin{enumerate}
    \item Placer le logo dans différents contextes (site web, affiche, carte de visite).
    \item Vérifier la lisibilité.
\end{enumerate}

\subsection{Les critères de lisibilité}

\subsubsection{La taille des détails}

Les détails doivent être visibles à petite taille. Si les détails disparaissent, le logo est trop complexe.

\subsubsection{Le contraste}

Le contraste doit être suffisant à toutes les tailles. Si le contraste est faible, le logo est illisible.

\subsubsection{La netteté}

La netteté doit être parfaite à toutes les tailles. Si les bords sont flous, le logo est pixellisé.

\subsubsection{La reconnaissance}

La reconnaissance doit être immédiate à toutes les tailles. Si le logo n'est pas reconnaissable, il est trop complexe.

\subsection{Les bonnes pratiques}

\begin{itemize}
    \item \terme{Tester} à toutes les tailles.
    \item \terme{Tester} à l'écran.
    \item \terme{Tester} sur papier.
    \item \terme{Tester} en contexte.
    \item \terme{Vérifier} les détails.
    \item \terme{Vérifier} le contraste.
    \item \terme{Vérifier} la netteté.
    \item \terme{Vérifier} la reconnaissance.
    \item \terme{Simplifier} si nécessaire.
\end{itemize}

\subsection*{Conclusion de la section}

La \terme{lisibilité} est essentielle pour un logo. Il faut tester le logo à différentes tailles : très petite, petite, moyenne, grande. On peut tester à l'écran, sur papier, en contexte. Les critères sont : taille des détails, contraste, netteté, reconnaissance. Il est essentiel de tester à toutes les tailles, de vérifier les détails, de simplifier si nécessaire.

Maintenant que nous savons évaluer la lisibilité, nous allons, dans la section suivante, apprendre à \terme{tester les versions monochromes}.

% ============================================================
\section{Tester les versions monochromes}
% ============================================================

\subsection{Pourquoi tester en noir et blanc ?}

Un logo doit fonctionner en \terme{noir et blanc}, pour plusieurs raisons :

\begin{itemize}
    \item \terme{Impression} : certains supports impriment en noir et blanc.
    \item \terme{Gravure} : certains supports gravent en noir et blanc.
    \item \terme{Fax} : certains documents sont faxés en noir et blanc.
    \item \terme{Photocopie} : certains documents sont photocopiés.
    \item \terme{Marque} : la version monochrome est souvent utilisée.
    \item \terme{Test} : la version monochrome révèle les défauts.
\end{itemize}

Un logo qui ne fonctionne pas en noir et blanc est un logo fragile.

\subsection{Les versions monochromes}

\subsubsection{Noir sur blanc}

La version \terme{noir sur blanc} est la plus courante. Elle doit être parfaite.

\subsubsection{Blanc sur noir}

La version \terme{blanc sur noir} est utilisée sur fond sombre. Elle doit être lisible.

\subsubsection{Niveaux de gris}

La version en \terme{niveaux de gris} est utilisée pour les impressions en noir et blanc. Elle doit conserver les contrastes.

\subsubsection{Monochrome couleur}

La version \terme{monochrome couleur} utilise une seule couleur (par exemple, le bleu de la marque). Elle doit fonctionner.

\subsection{Tester en noir et blanc}

\subsubsection{Test à l'écran}

\begin{enumerate}
    \item Convertir le logo en noir et blanc.
    \item Vérifier la lisibilité.
    \item Vérifier les contrastes.
\end{enumerate}

\subsubsection{Test sur papier}

\begin{enumerate}
    \item Imprimer le logo en noir et blanc.
    \item Vérifier la lisibilité.
    \item Vérifier la qualité.
\end{enumerate}

\subsubsection{Test en contexte}

\begin{enumerate}
    \item Placer le logo en noir et blanc dans différents contextes.
    \item Vérifier la lisibilité.
\end{enumerate}

\subsection{Les problèmes courants}

\subsubsection{Perte de détails}

Certains détails disparaissent en noir et blanc (couleurs similaires, contrastes faibles).

\subsubsection{Perte de reconnaissance}

Le logo n'est plus reconnaissable en noir et blanc (trop dépendant de la couleur).

\subsubsection{Perte d'équilibre}

Le logo est déséquilibré en noir et blanc (masses qui changent).

\subsection{Les bonnes pratiques}

\begin{itemize}
    \item \terme{Tester} en noir sur blanc.
    \item \terme{Tester} en blanc sur noir.
    \item \terme{Tester} en niveaux de gris.
    \item \terme{Tester} en monochrome couleur.
    \item \terme{Vérifier} la lisibilité.
    \item \terme{Vérifier} la reconnaissance.
    \item \terme{Vérifier} l'équilibre.
    \item \terme{Simplifier} si nécessaire.
\end{itemize}

\subsection*{Conclusion de la section}

Un logo doit fonctionner en \terme{noir et blanc}. Il faut tester le logo en noir sur blanc, en blanc sur noir, en niveaux de gris, en monochrome couleur. Les problèmes courants sont : perte de détails, perte de reconnaissance, perte d'équilibre. Il est essentiel de tester en noir et blanc, de vérifier la lisibilité, de simplifier si nécessaire.

Maintenant que nous savons tester les versions monochromes, nous allons, dans la section suivante, apprendre à \terme{préparer les variantes d'un logo}.

% ============================================================
\section{Préparer les variantes d'un logo}
% ============================================================

\subsection{Pourquoi préparer des variantes ?}

Un logo doit être \terme{adaptable} à différents contextes. Il faut donc préparer des \terme{variantes} :

\begin{itemize}
    \item \terme{Horizontale} : pour les espaces larges.
    \item \terme{Verticale} : pour les espaces étroits.
    \item \terme{Carrée} : pour les avatars, les favicons.
    \item \terme{Symbole seul} : pour les petites tailles.
    \item \terme{Texte seul} : pour les contextes où le symbole n'est pas nécessaire.
    \item \terme{Monochrome} : pour les impressions en noir et blanc.
\end{itemize}

Les variantes permettent au logo de fonctionner dans tous les contextes.

\subsection{Les variantes principales}

\subsubsection{Version horizontale}

La version \terme{horizontale} présente le symbole et le texte côte à côte. Elle est utilisée dans les espaces larges (en-tête de site, bannière, signature email).

\subsubsection{Version verticale}

La version \terme{verticale} présente le symbole au-dessus du texte. Elle est utilisée dans les espaces étroits (avatar, favicon, carte de visite carrée).

\subsubsection{Version carrée}

La version \terme{carrée} présente le logo dans un carré. Elle est utilisée pour les avatars, les favicons, les applications.

\subsubsection{Symbole seul}

Le \terme{symbole seul} est utilisé pour les petites tailles (favicon, avatar, icône). Il doit être reconnaissable.

\subsubsection{Texte seul}

Le \terme{texte seul} est utilisé pour les contextes où le symbole n'est pas nécessaire (signature, mention). Il doit être lisible.

\subsubsection{Monochrome}

La version \terme{monochrome} est utilisée pour les impressions en noir et blanc, les gravures, les fax.

\subsection{Préparer les variantes}

\subsubsection{Étape 1 : Créer la version principale}

Créer la version principale du logo, avec toutes les options.

\subsubsection{Étape 2 : Créer les variantes}

Créer les variantes en adaptant la composition :

\begin{itemize}
    \item \terme{Modifier} la position des éléments.
    \item \terme{Modifier} les proportions.
    \item \terme{Modifier} les espacements.
\end{itemize}

\subsubsection{Étape 3 : Vérifier}

Vérifier chaque variante :

\begin{itemize}
    \item \terme{Lisibilité} à toutes les tailles.
    \item \terme{Cohérence} avec la version principale.
    \item \terme{Reconnaissance}.
\end{itemize}

\subsubsection{Étape 4 : Organiser}

Organiser les variantes :

\begin{itemize}
    \item \terme{Nommer} chaque variante.
    \item \terme{Classer} par type.
    \item \terme{Documenter} les usages.
\end{itemize}

\subsection{Les bonnes pratiques}

\begin{itemize}
    \item \terme{Créer} plusieurs variantes.
    \item \terme{Vérifier} chaque variante.
    \item \terme{Maintenir} la cohérence.
    \item \terme{Nommer} les variantes.
    \item \terme{Documenter} les usages.
    \item \terme{Préparer} les fichiers pour chaque usage.
    \item \terme{Tester} à toutes les tailles.
\end{itemize}

\subsection*{Conclusion de la section}

Les \terme{variantes} d'un logo permettent de l'adapter à différents contextes : horizontale, verticale, carrée, symbole seul, texte seul, monochrome. La préparation suit plusieurs étapes : créer la version principale, créer les variantes, vérifier, organiser. Il est essentiel de créer plusieurs variantes, de vérifier chaque variante, de maintenir la cohérence.

Maintenant que nous savons préparer les variantes, nous allons, dans la dernière section de ce chapitre, apprendre à \terme{organiser les fichiers sources et les exports}.

% ============================================================
\section{Organiser les fichiers sources et les exports}
% ============================================================

\subsection{Pourquoi organiser les fichiers ?}

L'\terme{organisation} des fichiers est essentielle pour :

\begin{itemize}
    \item \terme{Retrouver} facilement les fichiers.
    \item \terme{Modifier} le logo ultérieurement.
    \item \terme{Partager} les fichiers avec le client.
    \item \terme{Collaborer} avec d'autres designers.
    \item \terme{Archiver} le projet.
    \item \terme{Assurer} la cohérence.
\end{itemize}

Un projet bien organisé est un projet professionnel.

\subsection{Les fichiers sources}

\subsubsection{Le fichier AI}

Le \terme{fichier AI} est le fichier source principal. Il contient :

\begin{itemize}
    \item \terme{Le} logo.
    \item \terme{Les} calques organisés.
    \item \terme{Les} variantes.
    \item \terme{Les} couleurs globales.
    \item \terme{Les} guides.
\end{itemize}

Il faut toujours conserver le fichier AI.

\subsubsection{Le fichier de variantes}

Il est recommandé de créer un \terme{fichier de variantes} séparé, contenant toutes les versions du logo.

\subsubsection{Le fichier de charte}

Il est recommandé de créer un \terme{fichier de charte} (ou une planche de présentation), contenant :

\begin{itemize}
    \item \terme{Les} variantes.
    \item \terme{Les} couleurs.
    \item \terme{Les} typographies.
    \item \terme{Les} règles d'utilisation.
    \item \terme{Les} exemples d'application.
\end{itemize}

\subsection{Les exports}

\subsubsection{Pour l'impression}

\begin{itemize}
    \item \terme{PDF} : pour l'impression, avec traits de coupe et fonds perdus.
    \item \terme{EPS} : pour l'impression, si l'imprimeur le demande.
    \item \terme{TIFF} : pour les images haute résolution.
\end{itemize}

\subsubsection{Pour le web}

\begin{itemize}
    \item \terme{SVG} : pour les logos vectoriels sur le web.
    \item \terme{PNG} : pour les logos avec transparence.
    \item \terme{JPEG} : pour les logos sans transparence.
\end{itemize}

\subsubsection{Pour les réseaux sociaux}

\begin{itemize}
    \item \terme{PNG} : pour les avatars, les images de couverture.
    \item \terme{JPEG} : pour les posts.
\end{itemize}

\subsubsection{Pour les applications}

\begin{itemize}
    \item \terme{PNG} : pour les icônes (différentes tailles).
    \item \terme{SVG} : pour les icônes vectorielles.
\end{itemize}

\subsection{Organiser les fichiers}

\subsubsection{La structure des dossiers}

\begin{verbatim}
projet-logo/
  |---- sources/
  |   |---- logo-principal.ai
  |   |---- logo-variantes.ai
  |   `-- logo-charte.ai
  |---- exports/
  |   |---- impression/
  |   |   |---- logo-pdf.pdf
  |   |   |---- logo-eps.eps
  |   |   `-- logo-tiff.tiff
  |   |---- web/
  |   |   |---- logo-svg.svg
  |   |   |---- logo-png.png
  |   |   `-- logo-jpeg.jpg
  |   |---- reseaux-sociaux/
  |   |   |---- logo-avatar.png
  |   |   `-- logo-couverture.png
  |   `-- applications/
  |       |---- logo-icon-16x16.png
  |       |---- logo-icon-32x32.png
  |       `-- logo-icon-64x64.png
  |---- charte/
  |   |---- charte-graphique.pdf
  |   `-- charte-graphique.ai
  `-- documentation/
      |---- brief.pdf
      `-- notes.md
\end{verbatim}

\subsubsection{Les conventions de nommage}

\begin{itemize}
    \item \terme{Nom du projet} : en premier.
    \item \terme{Type} : logo, icône, avatar.
    \item \terme{Variante} : horizontal, vertical, carré.
    \item \terme{Format} : ai, svg, png, pdf.
    \item \terme{Version} : v1, v2, etc.
\end{itemize}

\subsubsection{Les versions}

\begin{itemize}
    \item \terme{Version 1} : première version.
    \item \terme{Version 2} : corrections.
    \item \terme{Version finale} : version validée.
\end{itemize}

\subsection{Les bonnes pratiques}

\begin{itemize}
    \item \terme{Conserver} le fichier AI.
    \item \terme{Organiser} les dossiers.
    \item \terme{Nommer} les fichiers clairement.
    \item \terme{Documenter} les usages.
    \item \terme{Exporter} dans les bons formats.
    \item \terme{Vérifier} les fichiers.
    \item \terme{Archiver} le projet.
    \item \terme{Communiquer} avec le client.
\end{itemize}

\subsection*{Conclusion de la section}

L'\terme{organisation} des fichiers sources et des exports est essentielle pour retrouver, modifier, partager, collaborer, archiver. On conserve le fichier AI, on prépare les variantes, on crée une charte, on exporte dans les bons formats (PDF, EPS, TIFF pour l'impression ; SVG, PNG, JPEG pour le web ; PNG, JPEG pour les réseaux sociaux ; PNG, SVG pour les applications). Il est essentiel de conserver le fichier AI, d'organiser les dossiers, de nommer les fichiers, d'exporter dans les bons formats.

% ============================================================
\section*{Conclusion du chapitre}
\addcontentsline{toc}{section}{Conclusion du chapitre}
% ============================================================

Nous voici au terme du chapitre 20. Nous avons appris à maîtriser la création de logos professionnels dans Illustrator :

\begin{itemize}
    \item \terme{Fonction d'un logo} : identifier, différencier, communiquer, rassurer, fédérer, marquer ; qualités (simplicité, mémorisation, intemporalité, versatilité, pertinence, appropriation).
    \item \terme{Catégories de logos} : typographique, symbolique, mixte, abstrait, emblème, lettrage, monogramme.
    \item \terme{Analyse du brief} : contexte, projet, objectifs, public, contraintes, références ; questions à poser.
    \item \terme{Recherche visuelle} : sources, méthode, moodboard.
    \item \terme{Esquisse} : quantité, rapidité, diversité, simplicité ; étapes.
    \item \terme{Logo géométrique} : formes simples, précision, équilibre, cohérence ; construction.
    \item \terme{Logo typographique} : typographie, composition, personnalisation ; construction.
    \item \terme{Symbole vectoriel} : simplicité, reconnaissance, mémorisation, versatilité ; construction.
    \item \terme{Lisibilité} : tailles à tester, critères, bonnes pratiques.
    \item \terme{Versions monochromes} : noir sur blanc, blanc sur noir, niveaux de gris, monochrome couleur.
    \item \terme{Variantes} : horizontale, verticale, carrée, symbole seul, texte seul, monochrome.
    \item \terme{Organisation des fichiers} : sources, exports, charte, documentation.
\end{itemize}

Le logo est bien plus qu'une simple image : c'est la signature visuelle d'une marque, d'une entreprise, d'une organisation, d'un produit. Il permet d'identifier, de reconnaître, de mémoriser. Il transmet des valeurs, une personnalité, une promesse. Un bon logo est simple, mémorable, intemporel, versatile, pertinent, approprié. Un mauvais logo est complexe, daté, illisible, incohérent.

La création d'un logo est un processus stratégique, pas seulement esthétique. Elle repose sur une compréhension profonde du client, de son marché, de son public, de ses valeurs. Elle exige de la recherche, de l'analyse, de l'esquisse, de l'itération, de la précision. Un logo ne se crée pas en quelques minutes : il se construit, se travaille, se peaufine.

Ce chapitre t'a donné les bases de la création de logos professionnels. Mais la maîtrise vient avec la pratique. N'hésite pas à t'entraîner sur des projets variés : marques fictives, refontes de logos, exercices créatifs. Chaque projet est un nouveau défi, et chaque défi est une occasion de progresser.

Dans le chapitre suivant, nous aborderons l'\terme{illustration vectorielle et l'identité visuelle}. Tu apprendras à créer des icônes, à concevoir des pictogrammes, à créer des illustrations simples, à harmoniser les formes et les couleurs, à développer un système graphique cohérent, à concevoir une palette de marque, à associer les couleurs et les typographies, à créer des éléments de papeterie, et à présenter une identité visuelle. Ces compétences te permettront de créer des identités visuelles complètes.

\vspace{2em}

\begin{center}
\textit{Fin du chapitre 20 --- Création de logos professionnels.}

\vspace{1em}

\textbf{Prochain chapitre : Chapitre 21 --- Illustration vectorielle et identité visuelle.}
\end{center}